% Subsecciones iniciales de la sección 3 (de la entrega P1, condensadas). Rol 4.
\subsection{Fuente y descripción del dataset}
\label{sec:fuente}
Los datos provienen de SENAMHI y están publicados en la Plataforma Nacional de Datos Abiertos como ``Monitoreo de los contaminantes del aire en Lima Metropolitana''~\cite{senamhi_datos_2025}, con licencia Open Data Commons Attribution. Se usa el archivo filtrado a la estación Santa Anita, con 82.542 registros y 15 columnas entre enero de 2015 y mayo de 2024; se verificó que coincide con los registros de esa estación en el archivo completo de la red. Cada fila contiene identificador, estación, fecha y hora, ubicación, concentraciones de PM$_{10}$, PM$_{2{,}5}$ y NO$_2$ (columna \texttt{PM2\_5} para la variable de interés) y datos administrativos. Por tipo, la fecha y hora es temporal, las tres concentraciones son numéricas continuas, y el identificador, la estación y los datos administrativos son categóricos o de texto; la variable objetivo es PM$_{2{,}5}$ en la hora siguiente (regresión, por lo que no hay clases que balancear). El diccionario no indica la unidad; el informe de vigilancia de SENAMHI~\cite{senamhi_vigilancia_2021} reporta estos contaminantes en \textmu g/m\textsuperscript{3}, unidad que se adopta. Los identificadores, la ubicación y la fecha de corte no se usan como predictores, porque no describen el comportamiento del contaminante en una única estación.

\subsection{Selección de la estación y del periodo}
\label{sec:periodo}
Se trabaja con una sola estación para mantener una referencia geográfica constante; los resultados no representan a toda Lima. La disponibilidad de PM$_{2{,}5}$ varía mucho entre años: entre 84,6~\% y 94,3~\% en 2015--2017, y luego 50,2~\% en 2018, 20,9~\% en 2019, 0~\% en 2020 y 3,5~\% en 2021, hasta volver a 94,0~\%, 95,0~\% y 97,7~\% en 2022, 2023 y enero--mayo de 2024. Por ello se eligió el tramo continuo más reciente con alta disponibilidad, del 1 de enero de 2022 al 30 de abril de 2024, sin registros duplicados por fecha y hora. Comprende 20.424 horas, de las cuales 19.391 (94,94~\%) tienen medición de PM$_{2{,}5}$. Esta elección introduce un sesgo: el conjunto corresponde a los años en que la estación funcionó con normalidad y excluye el periodo de 2020, y dos años completos más cuatro meses caracterizan limitadamente la variación entre años.
